\documentclass[11pt,a4paper]{article}

\usepackage[utf8]{inputenc}
\usepackage[T1]{fontenc}
\usepackage[spanish,es-noquoting]{babel}
\usepackage[margin=2.5cm]{geometry}
\usepackage{amsmath}
\usepackage{booktabs}
\usepackage{hyperref}

\hypersetup{colorlinks=true, linkcolor=black, urlcolor=blue}

\title{Canal interno por la ecuación (3)\\
       \large Réplica del primer caso de estudio de O. Arsenyeva et al. (Tabla 5)}
\author{Informe de cambios en el código}
\date{16 de septiembre de 2026}

\begin{document}
\maketitle

\section{Encargo}

Reorientar \texttt{modelo\_f3.py} para que calcule \textbf{únicamente el canal interno} y
replique la Tabla 5 del artículo\footnote{O. Arsenyeva, J. Tran, M. Piper, E. Kenig,
\emph{An approach for pillow plate heat exchangers design for single-phase applications},
Appl. Therm. Eng. 147 (2019) 579--591.}, con las propiedades del agua del primer caso de
estudio, partiendo del número de chapas y la separación $b$ que da esa tabla, y \textbf{sin
imponer la pérdida de carga}.

\section{Cómo se ha hecho}

El fichero se reescribe por completo, conservando la estructura de
\texttt{modelo\_iterativo.py}: \emph{dataclass} de resultados, funciones de impresión, tabla de
datos de entrada y bucle de cálculo. Desaparece todo el bloque del canal externo y, con él, el
coeficiente global $U$.

La cadena de cálculo queda reducida a:
\begin{equation}
  w = \frac{G}{\rho\,N\,f_{chI}},\qquad
  Re = \frac{\rho\,w\,d_{eI}}{\mu},\qquad
  Nu = n_3 Re^{n_4} Pr^{n_5},\qquad
  h = \frac{Nu\,k}{d_{eI}},
\end{equation}
con $d_{eI}$ de la ecuación (14) y $f_{chI}$ de la (16). La pérdida de carga se obtiene
después, como \emph{resultado}, de la ecuación (19).

\subsection{Tres decisiones que conviene justificar}

\paragraph{La longitud $L_F$ se toma de la Tabla 5.} Al eliminar el canal externo no hay
coeficiente global $U$, y sin $U$ la ecuación (26) no puede dar la longitud. Como el encargo
pide explícitamente no fijar la pérdida de carga, hace falta una longitud de algún sitio para
poder evaluar la ecuación (19). Se toma la que reporta la propia Tabla 5 (2.072, 1.200 y
2.230\,m), y se usa \emph{solo} para la pérdida de carga: no interviene en la velocidad ni en
ninguno de los números adimensionales.

\paragraph{La separación $b$ no interviene.} Tanto $d_{eI} = 2\,b_i/\sqrt{2}$ (ec. 14) como
$f_{chI} = (b_i/\sqrt{2})(w_{pp}-2w_e)$ (ec. 16) dependen únicamente de $b_i$. La separación
entre paneles solo entra en la geometría del canal externo, que aquí no se calcula. Se
mantiene en la tabla de casos por completitud y porque así lo pedía el encargo, pero el
programa avisa por consola de que no afecta al resultado.

\paragraph{$n_{pl}$ son chapas.} Se conserva la constante \texttt{CHAPAS\_POR\_PANEL}, con la
que los canales internos son $n_{pl}/2$. El apartado de resultados recoge las dos lecturas.

\section{Datos de entrada}

Fluido frío (interno), de la Tabla 4: $\rho=980$\,kg/m$^3$, $\mu=8\cdot10^{-4}$\,Pa$\cdot$s,
$k=0.618$\,W/(m$\cdot$K), $c_p=4175$\,J/(kg$\cdot$K), de 10 a 50\,\textdegree C, con
$G=30$\,kg/s.

\begin{table}[h]
\centering
\begin{tabular}{lrrrr}
\toprule
Caso & $b_i$ [mm] & $n_{pl}$ & $b$ [mm] & $L_F$ [m]\\
\midrule
PPHE1 & 3.4 &  70 &  5.5 & 2.072\\
PPHE2 & 3.0 & 112 &  7.5 & 1.200\\
PPHE3 & 7.0 &  66 & 20.0 & 2.230\\
\bottomrule
\end{tabular}
\caption{$n_{pl}$, $b$ y $L_F$ de la Tabla 5; $b_i$ de la Tabla 1.}
\end{table}

\section{Resultados}

\begin{table}[h]
\centering
\small
\begin{tabular}{lrrrrrrrr}
\toprule
Caso & $N$ & $f_{chI}$ [m$^2$] & $d_{eI}$ [m] & $u$ [m/s] & $Re$ & $Nu$ & $\zeta$ & $h$ [W/m$^2$K]\\
\midrule
PPHE1 & 35 & 6.4912e-4 & 4.8083e-3 & 1.3474 & 7937 &  50.69 & 0.0794 & 6514.48\\
PPHE2 & 56 & 5.7276e-4 & 4.2426e-3 & 0.9544 & 4960 &  66.22 & 0.2449 & 9645.19\\
PPHE3 & 33 & 1.3364e-3 & 9.8995e-3 & 0.6941 & 8418 & 106.06 & 0.4162 & 6621.23\\
\bottomrule
\end{tabular}
\caption{Canal interno. $Pr = 5.4045$ en los tres casos.}
\end{table}

\subsection{Contraste con la Tabla 5 (T5)}

\begin{table}[h]
\centering
\begin{tabular}{lrrrrrrrrr}
\toprule
& \multicolumn{3}{c}{$u$ [m/s]} & \multicolumn{2}{c}{$Re$} & \multicolumn{2}{c}{$h$ [W/m$^2$K]} & \multicolumn{2}{c}{$\Delta P$ [kPa]}\\
\cmidrule(lr){2-4}\cmidrule(lr){5-6}\cmidrule(lr){7-8}\cmidrule(lr){9-10}
Caso & calc & T5 & cociente & calc & T5 & calc & T5 & calc & T5\\
\midrule
PPHE1 & 1.3474 & 1.835 & 1.362 & 7937 & 10810 & 6514 &  9804 & 33.09 & 60\\
PPHE2 & 0.9544 & 1.280 & 1.341 & 4960 &  6670 & 9645 & 11220 & 32.26 & 60\\
PPHE3 & 0.6941 & 0.940 & 1.354 & 8418 & 11390 & 6621 & 10206 & 22.84 & 60\\
\bottomrule
\end{tabular}
\caption{Lectura $N = n_{pl}/2$.}
\end{table}

\begin{table}[h]
\centering
\begin{tabular}{lrrrrrrr}
\toprule
& \multicolumn{3}{c}{$u$ [m/s]} & \multicolumn{2}{c}{$Re$} & \multicolumn{2}{c}{$\Delta P$ [kPa]}\\
\cmidrule(lr){2-4}\cmidrule(lr){5-6}\cmidrule(lr){7-8}
Caso & calc & T5 & cociente & calc & T5 & calc & T5\\
\midrule
PPHE1 & 0.6737 & 1.835 & 2.724 & 3968 & 10810 & 10.56 & 60\\
PPHE2 & 0.4772 & 1.280 & 2.682 & 2480 &  6670 &  9.22 & 60\\
PPHE3 & 0.3471 & 0.940 & 2.708 & 4209 & 11390 &  6.32 & 60\\
\bottomrule
\end{tabular}
\caption{Lectura $N = n_{pl}$.}
\end{table}

\section{Discusión}

\subsection{Ninguna de las dos lecturas reproduce la Tabla 5}

Con $N = n_{pl}/2$ la velocidad sale un 26\,\% por debajo; con $N = n_{pl}$, un 63\,\%. Lo
relevante no es la magnitud del error sino su \textbf{constancia}: los cocientes son 1.362,
1.341 y 1.354 en el primer caso, y 2.724, 2.682 y 2.708 en el segundo. Tres geometrías con
$b_i$ entre 3 y 7\,mm y $n_{pl}$ entre 66 y 112 dan el mismo factor dentro del 1.5\,\%.

Un error aleatorio no se comporta así. Un factor constante indica una diferencia
\emph{sistemática de definición} entre la sección de paso que usa el artículo y la de su
ecuación (16), o entre su recuento de canales y $n_{pl}$.

El programa lo cuantifica directamente con la columna $N$ implícito, es decir, los canales que
harían falta para dar la velocidad que reporta la Tabla 5:
\begin{equation}
  N_{\text{implícito}} = \frac{G}{\rho\,w_{T5}\,f_{chI}}
  = 25.70,\quad 41.76,\quad 24.37
\end{equation}
frente a los 70, 112 y 66 chapas de la tabla. El cociente $n_{pl}/N_{\text{implícito}}$ vale
2.72, 2.68 y 2.71.

\subsection{Coherencia con lo detectado anteriormente}

Este resultado concuerda con el análisis previo del primer caso de estudio, en el que ya se
determinó que las velocidades de la Tabla 5 implican un caudal frío de unos 81\,kg/s frente a
los 30\,kg/s de la Tabla 4, es decir, el mismo factor 2.72. El presente cálculo lo confirma por
una vía independiente y más directa, usando solo la ecuación (3) y sin intervenir en ningún
momento la ecuación (30) ni la constante discutida de la ecuación (28).

Conviene subrayar el contraste con el segundo caso de estudio: allí la lectura
$N = n_{pl}/2$ sí produce densidades coherentes y el área de la Tabla 9 se reproduce con un
error del 0.8\,\% en HE1. El problema es específico de la Tabla 5.

\subsection{Pérdida de carga}

Con la longitud de la propia Tabla 5, la pérdida de carga resulta de 33.09, 32.26 y
22.84\,kPa, frente a los 60\,kPa que el artículo declara agotados en las tres geometrías. El
déficit es coherente con la velocidad baja: al ser $\Delta P$ aproximadamente proporcional a
$w^2$, un factor 1.35 en velocidad son unos 1.8 en pérdida de carga, y en efecto
$60/33.09 = 1.81$.

Es decir, \textbf{el déficit de pérdida de carga no es un error adicional}, sino la misma
discrepancia de velocidad propagada. Esto es una comprobación interna útil: el modelo del canal
interno es consistente consigo mismo, y toda la desviación respecto al artículo entra por un
único punto, la velocidad.

\section{Ficheros}

\begin{table}[h]
\centering
\begin{tabular}{llp{7.8cm}}
\toprule
Fichero & Estado & Cambio\\
\midrule
\texttt{modelo\_f3.py} & Reescrito & Solo canal interno, caso de estudio 1, $n_{pl}$ y $b$ de la Tabla 5, $\Delta P$ como resultado\\
\bottomrule
\end{tabular}
\end{table}

No se ha modificado ningún otro fichero. \texttt{formulas.py},
\texttt{asignacion.py}, \texttt{geometría.py} e \texttt{iteracion\_vel.py} quedan intactos, y
\texttt{modelo\_iterativo.py} tampoco se toca.

\end{document}
